\section{Background and Related Works}

We focus on weight-only post-training quantization (PTQ) of LLMs in this work; other model-compression approaches include quantization-aware training (QAT) and pruning.
These methods are not strictly orthogonal to each other, as one could both prune and quantize a model.
Since \qts is a weight-only PTQ method, the rest of this section focuses on this area.
Most current \sota PTQ methods round to minimize the per-layer proxy loss from \citet{nagelround}.

\vspace{-0.15cm}
\begin{equation}
\label{eqn:proxy}
\ell(\hat W) = \Exv{x}{\|(\hat W - W)x\|^2} = \trace{(\hat W - W) H (\hat W - W)^T}
\end{equation}

Here, $\hat W \in \mathbb{R}^{m \times n}$ is the quantized weight matrix, $x \in \mathbb{R}^n$ is an input activation, and $H = \Exv{x}{xx^T} \in \mathbb{R}^{n\times n}$ is interpreted as a proxy Hessian matrix.
This objective is defined \textit{per-layer}, making it tractable for very large models.
However, minimizing it is difficult due to the non-differentiable nature of quantization.
Instead many works have proposed algorithms such as Hessian-based adaptive rounding, alternating optimization, and even coordinate descent to approximately minimize the proxy error \cite{aqlm,chee2023quip,tseng2024quip,optq}.

\subsection{Incoherence Processing}

The effectiveness of these methods depends on properties of $W$.
For example, many works have observed that weight and activation outliers cause poor quantization quality \cite{dettmers2022llmint8,awq,omniquant}.
In QuIP, \citet{chee2023quip} proposed that \textit{incoherence} was important for quantifying this effect.

\begin{definition}[\citet{chee2023quip}]
A Hessian $H \in \mathbb{R}^{n \times n}$ is $\mu$-incoherent if its eigendecomposition $H = Q \Lambda Q^T$ has $\max_{i,j} \; |Q_{ij}| = \max_{i,j} \; |e_i^T Q e_j| \leq \mu / \sqrt{n}$. 
A weight matrix $W \in \mathbb{R}^{m \times n}$ is $\mu$-incoherent if $\max_{i,j} \; \textstyle |W_{ij}| = \max_{i,j} \; |e_i^T W e_j| \leq \mu \|W\|_F / \sqrt{mn}$.
\end{definition}

Essentially, incoherence means the weights and important rounding directions (Hessian eigenvectors) are not too large in any direction, aiding quantization. 
To make $W, H$ incoherent (small $\mu$), one can perform \textit{incoherence processing} (IP) by conjugating $W, H$ with random orthogonal matrices $U, V: \tilde W \gets UWV^T, \tilde H \gets VHV^T$.
\qss introduced IP with the random Hadamard transformation (RHT), which performs $ \tilde W \gets V_mS_mWS_nV_n^T, \tilde H \gets V_nS_nHS_nV_n^T$ where $V_k$ is a  $k\times k$ Hadamard matrix and $S_k$ is a length $k$ random sign vector.
The RHT achieves, with probability $\ge 1 - \delta$, $\mu_{\tilde W} = 2 \log (4mn / \delta)$, meaning that $\tilde W$'s entries are approximately independently Gaussian distributed, which can aid quantization \cite{tseng2024quip,ashkboos2024quarot}.
We choose to build on incoherence processing here because the independent Gaussian-like weights it produces are suitable inputs for trellis coding~\citep{rptc}.

\subsection{Vector Quantization (VQ) for LLM PTQ}

$k$-bit VQ quantizes a $d$ dimensional vector $S$ to one of $2^{kd}$ $d$-dimensional vectors that form a codebook $C \in \mathbb{R}^{2^{kd} \times d}$ \cite{linde1980algorithm}.
Since $C$ is an unstructured collection of arbitrary vectors, VQ enables better shaping and packing density than scalar product quantization (SPQ), where each entry in $S$ is quantized independently \cite{6145679}. 
However, this also comes at the cost of exponential time quantization and exponential space inference: finding the nearest neighbor in $C$ requires $O(2^{kd}d)$ time, and storing $C$ requires $O(2^{kd}d)$ space.
%Even for a hypercube-shaped Uniform source distribution, VQ outperforms SPQ; SPQ is restricted to a uniform grid while VQ can use a lattice with better packing density (e.g. $A_2$ for $d=2$).
The current crop of \sota LLM PTQ methods, \qss and AQLM, both use VQ to achieve high-quality 2-bit models.
%Other methods such as GPTVQ have also applied VQ to improve quantization quality.
Since the shaping advantage of VQ comes from high dimensionality, both \qss and AQLM attempt to maximize dimensionality.
AQLM's uses a large 8D codebook (1MiB) that does not fit in L1 cache.
\qss uses an 8D compressible codebook based on the $E_8$ lattice, which is highly symmetric.
This codebook is compressible by $256\times$ and barely fits in L1 cache.
In either case, the VQ dimension is effectively hardware-limited to $\le 8$, motivating methods that enable even higher-dimensional quantization.
%This insight allows us to design fast codes for Gaussian sources, enabling fast inference.


% \subsection{Adaptive Rounding: How to Round}


% Instead, one way to quantize $W$ is by quantizing the columns of $W$ iteratively with linear feedback from \textit{already-quantized} weights.
% In this \textit{adaptive rounding} setup, $\hat W = Q(W + (W - \hat W)A)$, where $A \in \mathbb{R}^{n \times n}$ is a strictly $g$-block lower triangular feedback matrix and $Q$ performs $g$-dimensional quantization.
% QuIP and \qss showed that for $\mu$-incoherent $W, H$, adaptive rounding with $A = L_g - I$, where $L_gD_gL_g^T = H$ is the $g$-block LDL decomposition of $H$, was optimal among methods that specify $g$-block linear feedback as a function of $H$.
% This method, termed BlockLDLQ, achieves proxy error $\mu^2$rank($H$)$/n$ times lower than product quantization.

% Finally, while Equation \ref{eqn:proxy} formulates a per-layer proxy loss, it does not capture inter-layer interactions. 
% The same general concept of adaptive rounding can be applied at the model-level through fine-tuning. 
% For example, \qss fine-tunes layers to absorb the quantization error of already quantized layers before quantization, improving quantization quality with no inference-time cost.

\subsection{Trellis-Coded Quantization (TCQ)}

TCQ was first proposed by \citet{46532} to apply the benefits of trellis coded \textit{modulation}, a conceptually dual problem, to quantization.
Define a $(L, k, V)$ trellis $G$ as a directed graph with $2^L$ nodes, each of which has $2^{kV}$ incoming and outgoing edges and a value $\in \mathbb{R}^V$; these values form a codebook $C \in \mathbb{R}^{2^L\times V}$.
To quantize a length-$T$ sequence $S \in \mathbb{R}^T$, each contiguous length-$V$ subsequence of $S$ is assigned to a node $\in G$, with the restriction that the assigned nodes form a walk.
The reconstruction $\hat S$ of $S$ is then given by concatenating node values in the walk.
When $V = 1$, this setup describes \citet{46532}'s original scalar TCQ.
When $V > 1$, this describes TCVQ, which applies TCQ to vectors \cite{104316,156629}.

Finding the optimal $\hat S$ under an additive distortion metric can be done with the Viterbi algorithm in $O(2^LT)$ time.
This is linear in sequence length, enabling ultra-high dimensional quantization.
For exposition, we briefly describe the Viterbi algorithm here. Concretely, if we want to quantize a $T$-length scalar sequence reinterpreted as a sequence of vectors $s_1, s_2, \ldots, s_{T/V} \in \mathbb{R}^V$ using a trellis code with graph $G$ and codebook $C$, this corresponds to solving the optimization problem
\[
    \mbox{minimize} \; \sum_{i=1}^{T/V} \| C_{x_i} - s_i \|^2 \; \; \mbox{ over } x_1, x_2, \ldots, x_{T/V} \mbox{ the vertex sequence of a walk on graph } G.
\]
This optimization problem can be solved exactly with dynamic programming via the value function
\[
    \mathcal{V}_t(x) = \min \; \left\{ \sum_{i=1}^t \| C_{x_i} - s_i \|^2 \; \middle| \; x_1, x_2, \ldots, x_t \mbox{ the vertex sequence of a walk on } G \mbox{ and } x_t = x \right \}
\]
using the update rule
\[
    \mathcal{V}_t(y) = \min_{(x,y) \in G} \mathcal{V}_{t-1}(x) +  \| C_y - s_t \|^2.
\]
This Viterbi approach clearly takes time linear in $T$ and in the number of edges of $G$; with a few simple optimizations this can be brought to $O(2^L T)$. In comparison, brute-force-searching all possible $2^{kT}$ codes---which is what we would need to do for an unstructured $k$-bit $T$-dimensional codebook--- would take time proportional to $2^{LT/V}$. The ability to tractably find the closest representable vector in $\R^T$, even for large $T$, is in some sense the ``main benefit'' of trellis coding.
For i.i.d sources, as $L$ increases, TCQ efficiently approaches the infinite-length distortion-rate $D_R$, which lower bounds the attainable distortion of a $k$-bit quantizer \cite{6145679}.
As shown in Table~\ref{tab:distortion}, when quantizing an i.i.d.\ Gaussian with $k=2$, the scalar Lloyd-Max quantizer attains 0.118 MSE, \qs's 8D E8P codebook 0.089 MSE, our (\qt) 256D $L=16$ TCQ quantizer 0.069 MSE, and $D_R = 0.063$ \cite{lloyd1982least,1057548,tseng2024quip,10.5555/1146355}. 


% which helps because the finite block-length distortion-rate function (the information-theoretical lower bound of the attainable MSE) is decreasing in block length~\cite{6145679}.
% For example, when quantizing i.i.d. standard Gaussian at two bits per weight, a scalar Lloyd-Max quantizer~\cite{lloyd1982least,1057548} attains an MSE of $0.12$ with block length 1, the E8P codebook used by QuIP\# attains $0.09$ with block length 8, our TCQ-based quantizer attains $0.07$ with block length 256, while the information-theoretical lower bound with infinite block length is $0.0625$~\cite{10.5555/1146355}.

% In this setup, an allowable reconstructions $\hat S$ of $S$ length $T/V$ walks on $G$

% To quantize a sequence of $T$ weights $S \in \mathbb{R}^T$
% D
% These values form a $2^L \times V$ codebook to which $S$ can be quantized.
% Recall from section \ref{sec:intro} that , and that the optimal $\hat S$ for an additive distortion metric (e.g. squared error) can be found with the Viterbi algorithm in $O(2^LT)$ time. 
% Each group of $V$ weights in $S$ only needs to store which of the $2^{kV}$ edges it came from, giving an effective bitrate of $kV/V = k$.

% TCQ works by traversing a much large state space than $2^{kV}$ by taking advantage of the stateful decoding process.
% Note that, unlike VQ, TCQ performs \textit{better} when $V$ is small!


% \albert{Chris: write a related works section where we go through trellis cdoed quantization literature. Put why you should do TCQ in the intro.
% cite a bunch of info theory cited works
% too much notation
% recapitulate literature -- tcq was first proposed by blah, the idea here is that you have some blah, 
% finite state machine description
% qingyao write this as related works
% }

% In channel coding theory~\cite{6773024}, Trellis-coded modulation (TCM)~\cite{1056454,1093542} is a way to transmit information over a noisy channel reliably.
% Exploiting the duality of channel coding and source coding~\cite{shannon1959coding}, Trellis-coded quantization (TCQ) was proposed by \cite{46532} to quantize memoryless sources.
% The idea of TCQ is to create ``coupling'' among reconstructions by maintaining a finite-state machine.
% Due to its performance and modest complexity, TCQ is used in JPEG 2000~\cite{RABBANI20023}.
% Extensions of TCQ include TCVQ~\cite{104316,156629}, i.e. applying TCQ to vectors, which we also draw inspiration from.
% One distinctive feature of TCQ is its infinite block length, which helps because the finite block-length distortion-rate function (the information-theoretical lower bound of the attainable MSE) is decreasing in block length~\cite{6145679}.
% For example, when quantizing i.i.d. standard Gaussian at two bits per weight, a scalar Lloyd-Max quantizer~\cite{lloyd1982least,1057548} attains an MSE of $0.12$ with block length 1, the E8P codebook used by QuIP\# attains $0.09$ with block length 8, our TCQ-based quantizer attains $0.07$ with block length 256, while the information-theoretical lower bound with infinite block length is $0.0625$~\cite{10.5555/1146355}.
% In particular, we define the finite-length distortion-rate function $\mathbb{D}(k, n)$ as the lowest MSE for a quantizer that uses $k$ bits per sample and operates on $n$ samples.
% It can be shown that for i.i.d. standard Gaussian  $\mathbb{D}(k, n) = \min(1, 2^{-2(k-O(1/\sqrt{n}))})$~\cite{10.5555/1146355,6145679}.
% For example, if you quantize 256 weights with product quantization with group size eight at 2 bits per sample, the average distortion is at least $\mathbb{D}(2, 8) = ???$, but for TCQ the lower bound is $\mathbb{D}(2, 256) = ???$.


% Consider the problem of quantizing a random real length-$T$ sequence $S \in \mathbb{R}^T$ drawn from some source distribution $\mathcal{D}$ of sequences.
% $k$-bit quantization aims to find the $\hat S$ in a codebook $C \in \mathbb{R}^{2^{Tk} \times T}$ that minimizes some distortion metric $D(S, \hat S): \mathbb{R}^T \times C \to \mathbb{R}_{\ge 0}$.
% Here, we focus on additive distortion metrics, where $D(S, \hat S) = \sum_{i=1}^{T} d(S_i, \hat S_i)$ for some scalar distortion metric $d: \mathbb{R} \times \mathbb{R} \to \mathbb{R}_{\ge 0}$.
% In scalar quantization (SQ), each element $S_i \in \mathbb{R}$ of $S$ is quantized to a scalar codebook $c \in \mathbb{R}^{2^k}$ by $\hat S_i = \arg\min_{\hat S_i \in c} d(S_i, \hat S_i)$, giving $C = c^{T}$.
% SQ usually performs poorly due to suboptimal distribution shaping and packing density. 
% A natural extension is to perform vector quantization (VQ), where the entirety of $S$ is quantized to $C$ at once, i.e. $\hat S = \arg\min_{\hat S \in C} D(S, \hat S)$.
% This relaxes $C$ to contain arbitrary vectors and allows for improved shaping and packing.
% However, unstructured VQ costs exponential time and space in the bitrate $k$ and dimension $T$. 
% For large $k$ or $T$, this is intractable.

% Trellis coded quantization (TCQ) solves this intractability problem with a \textit{stateful} quantizer, where weights are quantized according to a trellis.
% Define a $(L, k, V)$ trellis $G$ as a simple directed graph with $2^L$ nodes, where each node has $2^{kV}$ incoming and outgoing edges and a value $\in \mathbb{R}^V$; these values form a codebook $P \in \mathbb{R}^{2^L \times V}$.
% In TCQ, $S$ is first split into $T/V$ contiguous $V$D vectors $S_{1 \le j \le T/V} \in \mathbb{R}^V$.
% Each $S_j$ is assigned to a node $G_j$ in $G$, with the restriction that $\exists$ an edge from $G_{j-1}$ to $G_j$.
% The set of representable length $T$ reconstructions from $G$ (i.e. $C$) is the set of all length $T/V$ walks on $G$, where node values are concatenated to give a reconstruction.
% Quantizing $S$ amounts to finding 
% %\hat S = [ \hat S_{1 \le j \le T/V} ]
% $\hat S = \arg\min_{\hat S_j \in P, \exists \mbox{ \scriptsize edge } \{G_{j-1}, G_j\} \forall j > 1} \sum_{j=1}^{T/V} d(S_j, p_j)$
% This can be done with the Viterbi algorithm, which takes $O(2^LT)$ time.
% Since each $\hat S_j$ only requires storing which outgoing edge of $G_{j-1}$ $\{G_{j-1}, G_j\}$ is, the effective bitrate is $Vk/V = k$, as desired.
% Figure \ref{fig:bitshift} shows an example of quantizing with a $(2, 1, 1)$ trellis.



% In TCQ, the codebook size $2^L \times V$ is independent of both the sequence length $T$ and bitrate $k$.
% TCQ can be interpreted as $T$D VQ to the codebook of all length $T/V$ walks on $G$, essentially enabling high dimensional vector quantization.
% During inference, TCQ requires knowing $G$ to reconstruct $\hat S$.
% A naive implementation requires storing all $2^{L + KV}V$ entries in $G$, which is expensive. 
% In \qt, we adopt the ``bitshift'' trellis defined by Mao et al.
% Specifically, node $i$ has edges to all nodes $j$ where the upper $L-kV$ bits of $\mbox{bin}(j)$ equal the lower $L-kV$ bits of $\mbox{bin}(i)$, where $\mbox{bin}(x)$ gives the $L$-length binary represention of $x$. 
% Here, neighboring weights share $L-kV$ bits and decoding the next weight only requires left-shifting $kV$ bits, a hardware-supported instruction on almost all platforms.



% Here, we describe TCQ for real numbers, but TCQ can be applied to any symbol space $\mathbb{S}$.
% In trellis coded quantization, a length-$T$ sequence $S \in \mathbb{R}^T$ is quantized accoriding a 

% \qingyao{Begin of Qingyao's description.}

% Consider a sequence of symbols $S = \{S_i \in \mathbb{S}: i \in [T]\}$ for some length $T \in \mathbb{N}_+$ and symbol space $\mathbb{S}$.
% In general, quantization consists of a pair of functions: a quantization function $A: \mathbb{S}^T \mapsto \left(\mathbb{F}^k\right)^T$ that maps $S$ to $Q = \{Q_i \in \mathbb{F}^{k}: i \in [T]\}$ for some rate $k \in \mathbb{R}_+$, and a dequantization function $B: \left(\mathbb{F}^k\right)^T \mapsto \mathbb{S}^T$ that maps $Q$ to $\hat{S} = \{\hat{S}_i \in \mathbb{S}: i \in [T]\}$.
% The goal of quantization is to minimize the distortion $D(S, \hat{S})$ for some distance measure $D: \mathbb{S}^T \times \mathbb{S}^T \to \mathbb{R}_+$.
% Here we focus on additive distance \qingyao{is this the right word?} of the form $D(S, \hat{S}) = \sum_{i=1}^T d(S_i, \hat{S}_i)$ for some element-wise distance measure $\mathbb{S} \times \mathbb{S} \to \mathbb{R}_+$.
% To describe TCQ, we first define $B$, and $A$ is implicitly defined as $A(S) = \arg\min_Q D(B(Q), S)$.

% Let $\delta_0$ be an initial state in some state space $\Theta$.
% We reconstruct one symbol per step: for each step, we update $\delta$ with $k$ bits from the input, i.e. $\delta_{i+1} = \Gamma(\delta_i, Q_i)$, where $\Gamma$ is a state transition function, also known as a ``trellis''.
% The reconstruction for the $i$-th symbol is $\hat{S}_i = C(\delta_i, Q_i)$, where $C: \Delta \times \mathbb{F}^k \mapsto \mathbb{S}$ is a codebook.
% % Note that the domain of $C$ does not depend on $T$.

% Note that $B$ is equivalent to the encoding process of a convolutional code, and $A(S) = \arg\min_Q D(B(Q), S)$ is equivalent to the maximum likelihood decoding of a convolutional code.
% To compute $A$, we use the Viterbi algorithm which takes $O(2^L T)$ time and $O(2^{L-k} T)$ space.
% We refer readers to XXX for a detailed treatment.
% In contrast, enumerating all possible $Q$ as in na\"ive VQ would take $O(2^{kT})$ time.

% \subsection{Random Permutation Trellis Coder (RPTC)}

% In particular, we adopt RPTC~\cite{5453439}, defined as follow: each symbol $S_i \in \mathbb{R}^V$ with $V = 2$, each state $\delta_i \in \Delta = \mathbb{F}^L \cong [0, 2^L - 1] \cap \mathbb{Z}$ is an $L$-bit shift register reinterpretable as an integer, the trellis $\Gamma(\delta_i, Q_i)$ shifts out the most significant $k$ bits of $\delta_i$ while shifts in the $k$ bits in $Q_i$, and the codebook $C(\delta_i, Q_i) = \Phi^{-1}\left(\Lambda(\delta_i)/2^L+1/2^{L+1}\right)$ where $\Phi$ is the standard Gaussian CDF and $\Lambda: \mathbb{F}^L \mapsto \mathbb{F}^L$ is a hash function.
% For good performance, $\Lambda$ must satisfy the condition that, if $Q$ consists of i.i.d. uniform random bits, then $B(Q)$ must resemble a random i.i.d. Gaussian sequence~\cite{5453439}.
% The identity function is a poor $\Lambda$ because it induces a strong correlation between adjacent $\hat{S}_i$.
% See Section~\ref{sec:qtip} for how we choose $\Lambda$ and $\delta_0$.

% We choose RPTC for two reasons.
% First, it is the \sota quantizer for i.i.d. samples, achieving an MSE of $\approx 0.07$ for $S_i \sim \mathcal{N}(0, 1)$ at 2 bits per sample, whereas E8P achieves $\approx 0.09$ and the information-theoretic lower bound predicted by the Gaussian distortion-rate function is $0.0625$.
% Recall that incoherence processing reduces LLMs to nearly i.i.d. Gaussian samples, so a lower Gaussian MSE suggests better performance of the quantized model.
% Second, $\delta_i$ only depends on $Q_i, Q_{i-1}, \cdots, Q_{i-\lceil L/k \rceil+1}$ as earlier bits are shifted out, which makes dequantization parallelizable, enabling fast inference.

% \qingyao{End of Qingyao's description.}



